\documentclass{szb-solution}

\area{Sorozatok}
\title{Numerikus sorozatok I}
\subject{Matematika G1}
\subjectCode{BMETE94BG01}
\date{Utoljára frissítve: \today}
\docno{4}

\begin{document}
\maketitle

\subsection{Konvergencia a definíció alapján}

\begin{enumerate}[a)]
  \item Legyen
        $$
          a_n=\left|\frac{n+1}{3n-8}\right|
          \text.
        $$
        A várható határérték $A=1/3$. Ha $n\geq3$, akkor $3n-8>0$, ezért
        $$
          \left|a_n-\frac13\right|
          =\left|\frac{n+1}{3n-8}-\frac13\right|
          =\frac{11}{9n-24}
          \text.
        $$
        Adott $\varepsilon>0$ esetén válasszuk az $N$ természetes számot úgy,
        hogy
        $$
          N>\max\left\{3,\frac{11}{9\varepsilon}+\frac83\right\}
          \text.
        $$
        Ekkor minden $n>N$ esetén
        $11/(9n-24)<\varepsilon$, tehát
        $$
          \boxed{a_n\longrightarrow\frac13}
          \text.
        $$

  \item A
        $$
          b_n=\frac{n(-1)^n-1}{2n}
          =\frac{(-1)^n}{2}-\frac{1}{2n}
        $$
        sorozat páros és páratlan indexű részsorozatai:
        $$
          b_{2k}=\frac12-\frac{1}{4k}\longrightarrow\frac12,
          \qquad
          b_{2k+1}=-\frac12-\frac{1}{4k+2}\longrightarrow-\frac12
          \text.
        $$
        A két részsorozat határértéke különböző, ezért
        $\boxed{(b_n)\text{ nem konvergens}}$.
\end{enumerate}

\subsection{Sorozatok határértéke}

\begin{enumerate}[a)]
  \item
        $$
          \lim_{n\to\infty}\frac{n^2-6n+7}{n^2+12n+49}
          =\lim_{n\to\infty}
          \frac{1-6/n+7/n^2}{1+12/n+49/n^2}
          =\boxed{1}
          \text.
        $$

  \item Mivel $1+2+\dots+n=n(n+1)/2$,
        \begin{align*}
          \frac{1+2+\dots+n}{n+2}-\frac n2
          &=\frac{n(n+1)-n(n+2)}{2(n+2)}
          =-\frac{n}{2n+4}
          \longrightarrow\boxed{-\frac12}
          \text.
        \end{align*}

  \item A számlálót és a nevezőt $3^n$-nel osztva
        $$
          \frac{(-2)^n+3^n}{(-2)^{n+1}+3^{n+1}}
          =\frac{(-2/3)^n+1}{-2(-2/3)^n+3}
          \longrightarrow\boxed{\frac13}
          \text.
        $$

  \item Az első tört számlálójában $n^{3/2}$, nevezőjében
        $\sqrt[4]{5}\,n^{3/2}$ a domináns rend. Ezért
        $$
          \frac{\sqrt{n^3+3n^2}+\sqrt[3]{n^4+1}}
          {\sqrt[4]{5n^6+2}+\sqrt[5]{n^7+3n^3}}
          \longrightarrow\frac{1}{\sqrt[4]{5}}
          \text.
        $$
        A második tagra
        $$
          0\leq\frac{n!}{(n+1)!+3^{2n}}
          \leq\frac{n!}{(n+1)!}=\frac{1}{n+1}\longrightarrow0
          \text,
        $$
        így a teljes határérték
        $$
          \boxed{\frac{1}{\sqrt[4]{5}}}
          \text.
        $$

  \item Konjugálttal bővítve
        \begin{align*}
          &\sqrt{n^2+n+1}-\sqrt{n^2-n+1}
          \\
          &\qquad
          =\frac{2n}{\sqrt{n^2+n+1}+\sqrt{n^2-n+1}}
          =\frac{2}{\sqrt{1+1/n+1/n^2}+\sqrt{1-1/n+1/n^2}}
          \longrightarrow\boxed{1}
          \text.
        \end{align*}

  \item Legyen $u_n=\sqrt[3]{n^2-n^3}$. A köbök összegének azonosságával
        $$
          u_n+n
          =\frac{u_n^3+n^3}{u_n^2-nu_n+n^2}
          =\frac{n^2}{u_n^2-nu_n+n^2}
          \text.
        $$
        Mivel $u_n=-n(1-1/n)^{1/3}$, ezért
        $$
          u_n+n
          =\frac{1}
          {(1-1/n)^{2/3}+(1-1/n)^{1/3}+1}
          \longrightarrow\boxed{\frac13}
          \text.
        $$
\end{enumerate}

\subsection{Rendőrelv}

Minden $n\geq1$ esetén $1\leq n\leq2^n$, ezért
$$
  0<\frac{1}{3^n}\leq\frac{n}{3^n}\leq\left(\frac23\right)^n
  \text.
$$
Mindkét szélső sorozat nullához tart, így a rendőrelv alapján
$$
  \boxed{\frac{n}{3^n}\longrightarrow0}
  \text.
$$

\subsection{További határértékek}

\begin{enumerate}[a)]
  \item
        \begin{align*}
          \sqrt[n]{5n^2-30n-21}
          &=n^{2/n}\left(5-\frac{30}{n}-\frac{21}{n^2}\right)^{1/n}
          \longrightarrow\boxed{1}
          \text.
        \end{align*}

  \item Mivel $|\cos n^3|\leq1$,
        $$
          \frac{\cos n^3}{2n}\longrightarrow0,
          \qquad
          \frac{3n}{6n+1}\longrightarrow\frac12
          \text,
        $$
        tehát
        $$
          \boxed{
            \lim_{n\to\infty}\left(
              \frac{\cos n^3}{2n}-\frac{3n}{6n+1}
            \right)=-\frac12
          }
          \text.
        $$

  \item Logaritmust véve
        $$
          n\ln\left(1+\frac{1}{n^2}\right)
          \sim\frac1n\longrightarrow0,
        $$
        ezért
        $$
          \boxed{\left(1+\frac{1}{n^2}\right)^n\longrightarrow1}
          \text.
        $$

  \item
        $$
          \left(\frac{3n-1}{3n+2}\right)^{2n}
          =\left(1-\frac{3}{3n+2}\right)^{2n}
          \longrightarrow
          e^{\lim\limits_{n\to\infty}-6n/(3n+2)}
          =\boxed{e^{-2}}
          \text.
        $$

  \item
        $$
          \ln\left(\left(1+\frac1n\right)^{\ln n}\right)
          =\ln n\cdot\ln\left(1+\frac1n\right)
          \sim\frac{\ln n}{n}\longrightarrow0
          \text,
        $$
        tehát
        $$
          \boxed{\left(1+\frac1n\right)^{\ln n}\longrightarrow1}
          \text.
        $$

  \item
        $$
          \left(\frac{n^2-n+1}{n^2+n+1}\right)^{2n+5}
          =\left(1-\frac{2n}{n^2+n+1}\right)^{2n+5}
          \text.
        $$
        Mivel
        $$
          (2n+5)\left(-\frac{2n}{n^2+n+1}\right)\longrightarrow-4,
        $$
        a határérték
        $$
          \boxed{e^{-4}}
          \text.
        $$
\end{enumerate}

\subsection{Az $e^k$ határérték teljes indukcióval}

Bizonyítsuk, hogy minden $k\geq0$ egészre
$$
  \lim_{n\to\infty}\left(1+\frac{k}{n}\right)^n=e^k
  \text.
$$
A $k=0$ eset nyilvánvaló, a $k=1$ eset pedig az $e$ szám sorozatos
definíciója. Tegyük fel, hogy az állítás valamely $k$-ra igaz. Ekkor
\begin{align*}
  \left(1+\frac{k+1}{n}\right)^n
  &=\left(1+\frac{k}{n+1}\right)^{n+1}
    \left(1+\frac1n\right)^n
    \frac{n+1}{n+k+1}
  \\
  &\longrightarrow e^k\cdot e\cdot1=e^{k+1}
  \text.
\end{align*}
Az indukciós lépés teljesült, tehát
$$
  \boxed{
    \lim_{n\to\infty}\left(1+\frac{k}{n}\right)^n=e^k
    \quad(k\in\mathbb N_0)
  }
  \text.
$$

\end{document}
